Usando os totais anuais da populacao mundial disponibilizados pelo World Bank para o indicador $\mathrm{SP.POP.TOTL}$, obtem-se:
\[
1970: 3{,}680{,}589{,}045,
\]
\[
1980: 4{,}437{,}602{,}892,
\]
\[
1990: 5{,}299{,}246{,}757,
\]
\[
2000: 6{,}161{,}884{,}811.
\]

Em bilhoes de habitantes, isso corresponde aproximadamente a
\[
3.68,\quad 4.44,\quad 5.30,\quad 6.16.
\]

Esses pontos mostram crescimento continuo da populacao mundial, mas com valores muito menores do que as extrapolacoes explosivas de muito longo prazo. Ao atualizar as Figuras 5.1 e 5.2, a curva deve passar por esses quatro pontos observados e manter a interpretacao de que o crescimento historico foi forte, porem bem mais moderado que uma projecao exponencial extrema sustentada por seculos.